\scanpage{121}
% archive foliation: 120
la condition (\unsure{83|82}) a) s'exprime par
\[
\tag{86}
\sum_i\varphi_i^\omega(\delta)=0
\]
la condition b) par la condition
\[
\tag{87}
k.((\varphi_i^\omega(\delta))_i)=F_t
\]
qui détermine d'emblée les $(\varphi_i^\omega(\delta))_i$ à scalaire inv. près --- la condition c les déterminant enfin (pour $\hat\imath$ également choisi) au \uline{signe} près.

Bien sûr, \deleted{on aura} \inserted{la relation entre $\kappa$ et les $\varphi_i^\omega(\delta)$ s'explicite par}
\[
\tag{88}
(\varphi_i^\omega(\delta))=
\underbrace{\check\kappa(\delta^\vee)}_{\mathcal O_t(-1)}
\wedge\omega
\in\mathcal O_t(-1)(\underline\omega)\simeq F_t\subset V_J(k)\subset k^J
\]
et la condition c) s'écrit Dieu sait comment --- il faudra \deleted{$\unsure{\varphi_i^\omega(\delta)}$} y revenir. Reste à exprimer la dépendance p.r. à $\omega$, pour \deleted{\unsure{exp}} trouver l'invariance de l'épinglage de Legendre p.r. au choix de $\hat\imath$ i.e. p.r. au choix de $\omega=\omega_{\hat\imath}$ ($\alpha:\mu_4^*(k)\simeq\underline\omega$ étant supposé déterminé), on trouve la relation nécessaire
\[
\tag{89}
\varphi_i^{\omega'}(\delta)=\deleted{\unsure{\imath_\omega^{-1}}}\imath_\omega\varphi_i^\omega(\delta)
\]
où
\[
\imath_\omega\overset{\text{déf}}=\alpha(\omega)
\qquad
\begin{gathered}
\alpha:\mu_4^*\simeq\underline\omega\ \text{étant l'iso.}\\
\text{(supposé déterminé)}\\
\text{associé à l'épinglage}\\
\text{de Legendre cherché}
\end{gathered}
\]

Posant
